\begin{onehalfspace}

\section{Il limite del calcolo classico}

Per decenni, la crescita dell'hardware classico descritta dalla Legge di Moore ha aumentato la potenza di calcolo. Questo progresso incontra però due limiti. Da un lato, i transistor si avvicinano alle dimensioni atomiche. Dall'altro, alcuni problemi richiedono risorse che crescono troppo rapidamente con la dimensione dell'input. È il caso della fattorizzazione di grandi interi, della ricerca non strutturata e della simulazione di sistemi quantistici: processori più potenti non bastano a cambiare questa crescita.

Il calcolo quantistico usa sovrapposizione, entanglement e interferenza per elaborare l'informazione. Feynman osservò che la simulazione di sistemi quantistici generici può richiedere risorse classiche esponenziali~\cite{feynman1982}. Deutsch definì poi il modello di calcolatore quantistico universale~\cite{deutsch1985}. Negli anni Novanta, Shor mostrò che una macchina quantistica ideale può risolvere la fattorizzazione e il logaritmo discreto in tempo polinomiale~\cite{shor1994}. Questi algoritmi mettono a rischio \ac{RSA} e altri sistemi a chiave pubblica. Un attacco concreto richiede però circuiti affidabili, molti qubit protetti e risorse per correggere gli errori.

Il rischio riguarda anche i dati cifrati oggi. Un avversario può conservarli e tentare di decifrarli in futuro. Sostituire i sistemi vulnerabili richiede inoltre anni di inventario, aggiornamento e verifica della compatibilità. Gli standard post-quantistici affrontano questo rischio sul piano crittografico. Sul piano del calcolo, occorre capire quali risorse servano a eseguire circuiti affidabili. Nell'era \ac{NISQ}, il numero di qubit da solo non basta: contano anche qualità delle porte, collegamenti fra qubit e accumulo degli errori.

Questa distanza fra potenziale algoritmico e realizzabilità costituisce il punto di partenza della tesi. Il contesto storico e fisico, dai qubit superconduttori alla rivoluzione algoritmica, è approfondito nelle Sezioni~\ref{app:intro:origini}--\ref{app:intro:algoritmi}; RSA, la transizione post-quantistica e il passaggio dall'era \ac{NISQ} ai sistemi fault-tolerant sono trattati nelle Sezioni~\ref{app:intro:rsa}--\ref{app:intro:nisq}. I fondamenti formali necessari all'analisi sono invece raccolti nella Sezione~\ref{chap:fondamenti}.

% -----------------------------------------------
\section{Motivazione e contributo del presente lavoro}

L'algoritmo di Shor è quindi un banco di prova particolarmente esigente: la sua complessità asintotica minaccia la crittografia moderna, mentre i circuiti profondi che ne realizzano il \textit{period finding} sono sensibili a decoerenza, errori di gate e misura. Il problema affrontato è comprendere quali interventi conservino informazione utile in presenza di rumore e, in particolare, \emph{dove} un modello appreso possa aggiungere valore rispetto a una regola analitica trasparente.

La tesi confronta due modi di intervenire. Il post-processing cerca candidati utili tra gli esiti già misurati, senza modificare lo stato durante il circuito. La correzione quantistica agisce invece durante il calcolo: usa più qubit e misura la sindrome, cioè un insieme di segnali che rivela gli errori senza leggere l'informazione logica. Questo richiede qubit, porte e calcolo classico aggiuntivi. Per valutare il beneficio, ogni tecnica viene confrontata con una regola di riferimento usando gli stessi dati e metriche coerenti.

Il lavoro procede in \textbf{due fasi}. La prima studia l'output di Shor su simulatore, con due configurazioni di rumore uniformi e illustrative. Queste non rappresentano la calibrazione di una specifica \ac{QPU}. Sugli stessi istogrammi si confrontano TOP-1 (Metodo~1), TOP-4 e il Metodo~2, che usa un classificatore \ac{ML} per decidere se applicare TOP-4. Riutilizzare gli stessi istogrammi evita che differenze casuali nella simulazione favoriscano una strategia. Il confronto con TOP-4 senza filtro, detto ablazione, isola poi il contributo del classificatore. Nei test il filtro non migliora significativamente la baseline ed è peggiore della sola TOP-4. Il beneficio viene quindi dalla ricerca di più candidati. Il protocollo e l'origine dei valori anticipati qui sono nella Sezione~\ref{app:intro:filo}; la verifica completa è nel Capitolo~\ref{chap:risultati}.

La \textbf{seconda fase}, nucleo della tesi, studia la correzione durante il calcolo. Si simulano tre codici di \ac{QEC}: ripetizione, Steane $[[7,1,3]]$ e surface code con decoder \ac{MWPM}. Per ciascuno si misura il tasso di errore logico $p_L$ al variare dell'errore fisico $p$ e della distanza $d$. Un'analisi separata introduce errori Pauli dopo le porte del circuito compilato di Shor. Questo modello serve a misurarne la sensibilità al rumore; il suo parametro per gate resta distinto dal tasso di fallimento della memoria misurato nei test QEC.

L'apprendimento automatico viene infine applicato alla \textbf{sindrome}, che contiene segnali diretti sugli errori. Sostituire \ac{MWPM} con una rete non dà un vantaggio stabile quando il modello di rumore è corretto. Affiancargli un decoder \textbf{ibrido} è invece utile in alcuni casi: la rete impara quando invertire la decisione del matching. Nei test con correlazioni assenti dal grafo del decoder analitico, l'errore logico si riduce fino al $21\%$. Quando il modello è adeguato, l'ibrido non mostra peggioramenti statisticamente rilevabili. Il risultato indica dove l'\ac{ML} può aiutare: nei casi in cui i dati contengono relazioni utili che il metodo di riferimento non rappresenta.

Il problema di ricerca, le sei domande, le ipotesi e il contratto sperimentale sono definiti nel Capitolo~\ref{chap:obiettivi}. Codice, sorgenti \LaTeX{}, seed e artefatti sono disponibili nel repository pubblico \url{https://github.com/claudio-dragotta/shor-nisq-qec-thesis}; la demo interattiva di Shor è all'indirizzo \url{https://shor-demo-6knp.onrender.com} (sorgenti: \url{https://github.com/claudio-dragotta/shor-demo}).

% -----------------------------------------------
\section{Struttura della tesi}

L'elaborato comprende sei capitoli e cinque appendici: dopo questa introduzione, il Capitolo~\ref{chap:obiettivi} formula problema e criteri; il Capitolo~\ref{chap:quadro} presenta teoria e metodologia; il Capitolo~\ref{chap:implementazione} descrive l'implementazione; il Capitolo~\ref{chap:risultati} espone e discute le due campagne, l'analisi di sensibilità e la QFT approssimata; il Capitolo~\ref{chap:conclusioni} risponde alle domande di ricerca. Le Appendici~\ref{app:teoria}--\ref{app:fase2} raccolgono, rispettivamente, teoria estesa, strumenti e codice, contesto e obiettivi, dettagli della fase classica e della fase QEC. La roadmap completa, con il contenuto referenziabile di ciascun capitolo e appendice, è trasferita nella Sezione~\ref{app:intro:struttura}.

\end{onehalfspace}
